\documentclass[11pt, oneside]{article} 
\usepackage{geometry}
\geometry{letterpaper} 
\usepackage{graphicx}
	
\usepackage{amssymb}
\usepackage{amsmath}
\usepackage{parskip}
\usepackage{color}
\usepackage{hyperref}

\graphicspath{{/Users/telliott/Github-math/figures/}}
% \begin{center} \includegraphics [scale=0.4] {gauss3.png} \end{center}

\title{Introduction}
\date{}

\begin{document}
\maketitle
\Large

%[my-super-duper-separator]

Let us summarize what we know so far about the basic theorems of geometry.  

They are all that you should really need to memorize to attack any of the other problems in the book and they'll be used repeatedly.

The first two aren't technically theorems but fundamental assumptions that we make about the geometrical world.

$\circ$ \ \ \hyperref[sec:supplementary_angle_theorem]{\textbf{supplementary angles}}

$\circ$ \ \  \hyperref[sec:alternate_interior_angle_theorem]{\textbf{alternate interior angles (parallel postulate)}}

The next three easily follow from those initial ideas.

$\circ$ \ \ \hyperref[sec:vertical_angle_theorem]{\textbf{vertical angles}}

$\circ$ \ \  \hyperref[sec:triangle_sum_theorem]{\textbf{triangle sum of angles}}

$\circ$ \ \ \hyperref[sec:complementary_angle_theorem]{\textbf{complementary angles}}

We have two basic methods for proving that two triangles are congruent

$\circ$ \ \ \hyperref[sec:SAS]{\textbf{SAS for congruence}}

$\circ$ \ \ \hyperref[sec:SSA_in_right]{\textbf{hypotenuse leg in a right triangle (HL)}}

(We proved that SSS is equivalent to SAS, and AAS is equivalent to ASA)

$\circ$ \ \ \hyperref[sec:SSS_implies_SAS]{\textbf{SSS implies SAS}}

Next we have the powerful fundamental theorems of geometry:

$\circ$ \ \  \hyperref[sec:isosceles_triangle_theorem]{\textbf{isosceles triangle theorem}} (sides $\rightarrow$ angles) and \hyperref[sec:isosceles_converse]{\textbf{converse}}

$\circ$ \ \ \hyperref[sec:external_angle_theorem]{\textbf{external angle theorem}}

$\circ$ \ \ \hyperref[sec:Thales_theorem]{\textbf{Thales' circle theorem}} (right angle in a semi-circle)

$\circ$ \ \ \hyperref[sec:area_ratio_theorem]{\textbf{area ratio theorem}}

$\circ$ \ \  \hyperref[sec:similar_right_triangles]{\textbf{similar triangles}} (similar $\triangle$s, same ratio of sides)

$\circ$ \ \ \hyperref[sec:Pythagoras_similar_triangles]{\textbf{Pythagorean theorem}} (similar triangles)

Some people might add a few more.  But I think this is a reasonable number to start with.  You should have these instantly available (and it's nice to know how to prove them, as well).

On similar triangles, so far we have only proved it for right triangles.  We have quite a bit more work to do.

Looking forward, probably the most important we have yet to do is

$\circ$ \ \ \hyperref[sec:peripheral_angle]{\textbf{inscribed angle theorem}} (on a circle is one-half central angle)

and the tremendously important general theorem that similar triangles have sides in the same ratio:

$\circ$ \ \ \hyperref[sec:similarity_and_ratios]{\textbf{similarity theorems}}

I've decided that we will sneak this proof in here.  If it's too difficult, it is OK to move on.  

\subsection*{Prop. I.5}

One of the most famous proofs in Euclid's book I is the forward version of the isosceles triangle theorem:  if two sides in a triangle are equal, then so are the opposing angles.  It is really only the second proof in all of Euclid (the first three propositions are constructions, and no. 4 is SAS).  It's difficult because Euclid proves it \emph{before} we know how to bisect an angle.

\label{sec:Euclid5}

\begin{center} \includegraphics [scale=0.5] {Euclid_1_5.png} \end{center}

Given AB = AC.  We will prove that the base angles are equal:  $\angle ABC = \angle ACB$.

\emph{Proof}.

Extend those two line segments so that AD = AE.  So then BD = CE by subtraction.  Connect BE and CD.

$\triangle ACD \cong \triangle ABE$ by SAS ($AB = AC$, $AD = AE$, and a shared angle at vertex $A$).

Therefore, the pair of marked angles in the left panel $\angle ABE = \angle ACD$.  And for the same reason, CD = BE. 

We also have that $BD = CE$, and the angle at D equals the angle at E.

Therefore, $\triangle BCD \cong \triangle BCE$ by SAS.

Thus, the angles marked with magenta dots are equal.

But now, black + magenta = red.  Since the latter two pairs are equal, by subtraction, so is the first.

$\angle ABC = \angle ACB$.

$\square$

Notice how the original figure is extended to provide auxiliary shapes helpful in the proof.  This is a common theme.

\subsection*{Prop. I.6}

We proved the converse of I.5 early in the book based on angle bisection.  Later, we will give another proof by contradiction.

The theorem together with its converse says that, in an isosceles triangle, the base angles are equal $\iff$ the two sides sides are equal (not the base).  

The symbol $\iff$ means \emph{if and only if}, so $A \iff B$ means that both $A \rightarrow B$ and $B \rightarrow A$.


\end{document}